% Options for packages loaded elsewhere
\PassOptionsToPackage{unicode}{hyperref}
\PassOptionsToPackage{hyphens}{url}
\PassOptionsToPackage{dvipsnames,svgnames,x11names}{xcolor}
%
\documentclass[
  11pt,
]{article}
\usepackage{amsmath,amssymb}
\usepackage{iftex}
\ifPDFTeX
  \usepackage[T1]{fontenc}
  \usepackage[utf8]{inputenc}
  \usepackage{textcomp} % provide euro and other symbols
\else % if luatex or xetex
  \usepackage{unicode-math} % this also loads fontspec
  \defaultfontfeatures{Scale=MatchLowercase}
  \defaultfontfeatures[\rmfamily]{Ligatures=TeX,Scale=1}
\fi
\usepackage{lmodern}
\ifPDFTeX\else
  % xetex/luatex font selection
\fi
% Use upquote if available, for straight quotes in verbatim environments
\IfFileExists{upquote.sty}{\usepackage{upquote}}{}
\IfFileExists{microtype.sty}{% use microtype if available
  \usepackage[]{microtype}
  \UseMicrotypeSet[protrusion]{basicmath} % disable protrusion for tt fonts
}{}
\makeatletter
\@ifundefined{KOMAClassName}{% if non-KOMA class
  \IfFileExists{parskip.sty}{%
    \usepackage{parskip}
  }{% else
    \setlength{\parindent}{0pt}
    \setlength{\parskip}{6pt plus 2pt minus 1pt}}
}{% if KOMA class
  \KOMAoptions{parskip=half}}
\makeatother
\usepackage{xcolor}
\usepackage[margin=1in]{geometry}
\setlength{\emergencystretch}{3em} % prevent overfull lines
\providecommand{\tightlist}{%
  \setlength{\itemsep}{0pt}\setlength{\parskip}{0pt}}
\setcounter{secnumdepth}{-\maxdimen} % remove section numbering
\ifLuaTeX
  \usepackage{selnolig}  % disable illegal ligatures
\fi
\usepackage{bookmark}
\IfFileExists{xurl.sty}{\usepackage{xurl}}{} % add URL line breaks if available
\urlstyle{same}
\hypersetup{
  pdftitle={Exact Independent Mutual-Visibility Chromatic Numbers of Strong Grid Graphs},
  pdfauthor={Ryutaro Yonezu --- Independent Researcher},
  colorlinks=true,
  linkcolor={Maroon},
  filecolor={Maroon},
  citecolor={Blue},
  urlcolor={blue},
  pdfcreator={LaTeX via pandoc}}

\title{Exact Independent Mutual-Visibility Chromatic Numbers of Strong
Grid Graphs}
\author{Ryutaro Yonezu --- Independent Researcher}
\date{18 November 2026}

\begin{document}
\maketitle

\textbf{Research No.73 · v1.0 · Unrefereed preprint.} Independent AI
audit: PASS (author-provided report); author Windows verification: PASS
(author-provided report). Human peer review and formal proof-assistant
verification: NOT PERFORMED. NOVELTY INCONCLUSIVE; REFERENCES PARTIAL.

\subsection{Abstract}\label{abstract}

For integers \(t\ge r\ge3\), we establish the exact value
\(\chi_{\mu_i}(P_t\boxtimes P_r)=\max\{4,\lceil r/2\rceil\}\) by an
explicit column coloring. Brešar, Peterin, Samadi and Yero (2026)
established the lower bound and equality for widths divisible by four
and widths three through seven. Their even-palette construction is
retained here. An odd-palette construction together with a general
geodesic lemma handles the remaining widths. The analytical argument
does not depend on finite computation. Historical novelty is
inconclusive; the independent AI audit reports a mathematical PASS
within its stated scope. Human peer review and formal proof-assistant
verification have not been performed.

\subsection{1. Definitions and lower
bound}\label{definitions-and-lower-bound}

Let \(G_{t,r}=P_t\boxtimes P_r\) have vertices \((i,j)\), \(0\le i<t\),
\(0\le j<r\). Distinct vertices are adjacent if the maximum of the two
absolute coordinate differences is one. The distance between \((i,j)\)
and \((I,J)\) is \(\max(|i-I|,|j-J|)\): each step changes either
coordinate by at most one, and taking diagonal steps followed by
axis-parallel steps attains this bound inside the rectangle.

An independent mutual-visibility (IMV) set \(S\) is independent and each
pair of distinct vertices in \(S\) admits a shortest path with no
internal vertex in \(S\). The IMV chromatic number is the minimum number
of such sets partitioning the vertex set.

\textbf{Theorem 1 (main claim).} For every integer \(t\ge r\ge3\),

\[\chi_{\mu_i}(G_{t,r})=\max\{4,\lceil r/2\rceil\}.\]

The lower bound is known from {[}1, Theorem 27{]}. For completeness, a
\(2\times2\) subgrid is a clique of size four. The diagonal
\(\{(j,j):0\le j<r\}\) has at most two vertices of any single IMV color,
because the unique geodesic joining any two of its vertices uses all
intermediate diagonal vertices: for endpoints \((a,a)\) and \((b,b)\),
\(a<b\), each of the \(b-a\) steps must increase both coordinates by
one. Three vertices of one color would therefore violate mutual
visibility. Counting diagonal vertices proves the other lower bound
\(\lceil r/2\rceil\).

\subsection{2. A column-phase lemma}\label{a-column-phase-lemma}

Write \(p(i,j)=(i+j)\bmod2\). Suppose \(t,r\ge2\) and a subset \(S\) of
\(G_{t,r}\) has the following properties:

\begin{enumerate}
\def\labelenumi{(\roman{enumi})}
\item
  the columns containing members of \(S\) are pairwise nonconsecutive;
\item
  in every such column \(j\), membership in \(S\) is precisely
  \(p(i,j)=p_j\), for one fixed bit \(p_j\);
\item
  the phase word of occupied columns in increasing order is a
  subsequence (not necessarily contiguous) of \((A,B,B,A)\), where
  \(A,B\in\{0,1\}\) and \(A\ne B\).
\end{enumerate}

\textbf{Lemma 2.} Under (i)--(iii), \(S\) is an IMV set.

Proof. Independence follows from (i) and (ii). For any two distinct
vertices of \(S\), write their absolute row and column differences as
\(h,d\).

If \(d=0\), then \(h\ge2\) by independence. Choose an adjacent column
within the grid; it is unoccupied by (i). Step diagonally into that
column toward the target row, walk vertically until one row before the
target, then step diagonally into the target. This path has exactly
\(h\) steps, and all its internal vertices are outside \(S\).

If \(d>0\), then \(d\ge2\). Condition (iii) implies that the endpoints
may be oriented so that the phase of the starting vertex differs from
the phase of every occupied column strictly between them. To see this,
examine their two indices in the parent four-letter word: consecutive
indices have no intervening letter; indices (1,3) and (2,4) have
intervening letter B and an endpoint A; indices (1,4) have two
intervening B's and endpoints A. Restricting to a subsequence cannot
create any new intervening letter. Reverse the endpoints if needed.
Write the resulting start as \((i,j)\) and target as \((I,J)\); set
\(u=\operatorname{sgn}(I-i)\), \(v=\operatorname{sgn}(J-j)\) (with
\(\operatorname{sgn}(0)=0\)). Every diagonal step preserves \(p\). Thus
any diagonal segment from the start through strictly intervening columns
avoids \(S\).

If \(h\ge d\), take \(d-1\) steps of displacement \((u,v)\), then
\(h-d\) steps of displacement \((u,0)\), then one step \((u,v)\) to the
target. The penultimate column \(J-v\) contains no member of \(S\) by
(i) because \(J\) is occupied. The initial diagonal segment avoids \(S\)
by the orientation above; the vertical segment remains in the unoccupied
column. All coordinates stay between the respective endpoint
coordinates. The path has length \(h\).

If \(h<d\), start with \(h\) steps \((u,v)\) and follow with \(d-h\)
horizontal steps \((0,v)\). If \(h=0\), the initial segment consists of
the start alone. Every internal vertex of the diagonal segment avoids
\(S\). For each internal vertex on the horizontal tail that lies in
\(S\) (thus strictly after step \(h\) and before step \(d\)), replace
its row \(I\) by an adjacent row inside the grid: \(I+1\) if \(I+1<t\),
otherwise \(I-1\). This flips its phase, hence removes it from \(S\) by
(ii). The changed columns are pairwise nonconsecutive by (i). Each
changed vertex therefore has its preceding and following vertices
unchanged and in row \(I\); this includes the interface vertex at step
\(h\) and the final endpoint. Their row differences are at most one and
column differences exactly one. No length or adjacency is lost. The
resulting \(d\)-step path stays inside the grid and avoids \(S\)
internally. It is a geodesic. This proves the lemma. QED.

\subsection{3. Explicit constructions}\label{explicit-constructions}

Set \(q=\max\{4,\lceil r/2\rceil\}\), so \(r\le2q\). Define a coloring
by a formula on the columns \(0,\dots,2q-1\), then retain only columns
\(0,\dots,r-1\) and all \(t\) rows. The proof will apply the lemma
directly to the smaller rectangle; it does not assume arbitrary
rectangular restrictions preserve mutual visibility.

\subsubsection{3.1 Odd palette}\label{odd-palette}

Suppose \(q\) is odd, so \(q\ge5\). Use the colors \(0,\dots,q-1\) and
put

\[c(i,j)=\bigl(2j+((i+j+\lfloor j/q\rfloor)\bmod2)\bigr)\bmod q.\]

A column has exactly the two possible colors \(\{2j,2j+1\}\pmod q\),
alternating with the row. Two consecutive columns have disjoint color
pairs: their four residues are consecutive modulo \(q\) and \(q\ge5\).
This verifies (i), and the alternating rows verify (ii).

Fix a color \(z\). Multiplication by two is invertible modulo odd \(q\),
so there are distinct representatives \(a,b\in\{0,\dots,q-1\}\) with
\(2a\equiv z\) and \(2b+1\equiv z\pmod q\). In the full \(2q\) columns,
precisely \(a,b,a+q,b+q\) contain \(z\). At \(a\) the phase is zero and
at \(b\) it is one; at \(a+q\) the phase is one and at \(b+q\) it is
zero, because the block flag \(\lfloor j/q\rfloor\) changes from zero to
one. If \(a<b\) the ordered word is \(0,1,1,0\); if \(b<a\) it is
\(1,0,0,1\). In either case (iii) holds. Removing a suffix of columns
preserves (i)--(iii). By Lemma 2, each color class is an IMV set.

\subsubsection{3.2 Even palette (known
pattern)}\label{even-palette-known-pattern}

If \(q=2k\) is even, \(k\ge2\). Use

\[c(i,j)=2(j\bmod k)+((i+j+\epsilon_j)\bmod2),\]

where \(\epsilon_j=1\) for \(k\le j<3k\) and zero otherwise. For each
\(a\in\{0,\dots,k-1\}\) the colors \(2a,2a+1\) occupy the four columns
\(a,a+k,a+2k,a+3k\) in the full width. Their phases are respectively
\(0,1,1,0\) and \(1,0,0,1\). These columns are separated by at least
two. Alternating row membership and retained column prefixes meet the
lemma's hypotheses. This formula restates the known even-palette pattern
from {[}1, Theorem 28{]}, with the present lemma establishing the
restriction directly.

Both cases yield an IMV coloring using at most \(q\) colors. Combined
with the lower bound, this proves Theorem 1. QED.

\subsection{4. Claim boundaries and
context}\label{claim-boundaries-and-context}

The lower bound, the even-palette full-width construction, and its phase
organization originate in {[}1{]}. The proposed contribution is the
odd-palette construction and a common geodesic argument for all widths,
including the justified restriction of the even construction. {[}1,
Theorem 28{]} explicitly covers \(t\ge4k\ge8\) and
\(t\ge r\in\{3,4,5,6,7\}\). The remaining widths in the stated range are
\(r\ge9\) with \(r\not\equiv0\pmod4\); the odd-palette widths
\(4k+1,4k+2\) (\(k\ge2\)) are the main target. Width \(4k+3\)
(\(k\ge2\)) is handled by the restriction argument for an even palette,
also when \(t=r\).

{[}2, Sections 1 and 6{]} explicitly separates independent
mutual-visibility coloring from total, outer and dual mutual-visibility
coloring and focuses on the latter parameters for strong grids. Its
equation (8) concerns ordinary MV coloring, without the independence
requirement. {[}3{]} studies a distance bound on mutual visibility, not
independence. These comparisons do not establish historical priority.
The literature search is targeted, not exhaustive (REFERENCES PARTIAL;
NOVELTY INCONCLUSIVE).

\subsection{5. Computational evidence (not used in
proof)}\label{computational-evidence-not-used-in-proof}

The accompanying standard-library Python scripts test the construction.
\texttt{verify\_paths.py} checks every same-color pair in 78 grids
(widths 3 through 28; heights r, r+1 and r+5), including independence,
path length, adjacency, boundaries and internal color avoidance.
\texttt{verify\_bfs.py} reconstructs the row-color table separately and
performs a graph BFS from each source, recording distances to same-color
terminal vertices without expanding them. This checks every same-color
pair in 40 grids (widths 3 through 22; heights r and r+3), plus proper
coloring along every edge. The implementations are separate but
developed in the same session and are not an external independent audit.
Local Linux reruns passed for 573,778 and 145,194 pairs respectively
(overlapping domains; do not add as a count of distinct pairs). No
finite test substitutes for the all-parameter proof. The author reports
Windows PowerShell PASS for the frozen v0.1 package. The author-provided
independent AI audit reports MODEL, LOWER, LEMMA, ODD, EVEN,
RESTRICTION, THEOREM, COMPUTATION and PACKAGE PASS. External C++ source
and raw logs were not supplied and were not independently rerun here.

\subsection{References}\label{references}

{[}1{]} B. Brešar, I. Peterin, B. Samadi, I. G. Yero, \emph{Independent
mutual-visibility coloring and related concepts}, Aequationes
Mathematicae 100 (2026), article 48. DOI:
https://doi.org/10.1007/s00010-026-01280-y . Theorems 27--28, publisher
HTML inspected 2026-11-18. Preprint: https://arxiv.org/abs/2505.04144 .

{[}2{]} S. Babu, M. Jakovac, D. Kuziak, A. Lakshmanan S., I. G. Yero,
\emph{A variety of the mutual-visibility coloring problem for graphs},
arXiv:2609.31427 (2026-09-25). https://arxiv.org/abs/2609.31427 .
Sections 1, 3 and 6 inspected in retrieved text, 2026-11-18.

{[}3{]} S. Babu, B. Brešar, A. Lakshmanan S., B. Samadi, \emph{Distance
mutual-visibility coloring: relations with (total) domination, exact
distance graphs and graph products}, arXiv:2510.10284 (2025).
https://arxiv.org/abs/2510.10284 . Abstract and Sections 1 and 5
inspected in the original arXiv HTML, 2026-11-18; broader literature
coverage remains partial.

\begin{center}\rule{0.5\linewidth}{0.5pt}\end{center}

AI-assisted preprint. Independent AI audit: PASS (author-provided scoped
report). Windows: PASS (author-provided report for frozen v0.1). Human
peer review: NOT PERFORMED. Formal proof-assistant verification: NOT
PERFORMED. NOVELTY INCONCLUSIVE; REFERENCES PARTIAL. Copyright (c) 2026
Ryutaro Yonezu. All rights reserved.

\end{document}
